\section{Introduction}\label{sec:intro}
%{\it Introduction:}\ 
Parton distribution functions (PDFs) provide an effective description of quarks and gluons inside a light-travelling nucleon~\cite{Ellis:1991qj,Thomas:2001kw}.
They play an essential role in calculating hardronic cross sections involving the nucleons~\cite{Gao:2017yyd}, and their uncertainties from phenomenological extractions have been one of the major
sources of errors in the theoretical predictions at Large Hadron Collider (LHC) and other hadron facilities. Thus, a precise knowledge of PDFs is very important both for the accurate tests of the Standard Model (SM) and for the search of new physics beyond the SM. On the other hand, the densities of partons in the nucleon provide direct information on its intrinsic properties such as the origin of the nucleon spin and mass, as well as the role of sea quarks for various physical quantities~\cite{Accardi:2012qut}. 
%Beyond PDF, there are also kinds of the functions proposed %to investigate the 3D picture of nucleon, such as %transverse momentum dependent PDF, generalized PDF and so on. 
  
However, directly calculating parton physics from first principles of quantum chromodynamics (QCD) has been a difficult task. A review on various efforts of doing so can be found in Ref.~\cite{Cichy:2018mum}. The proposal of the large-momentum effective field theory (LaMET) has been an important step toward meeting the challenge ~\cite{Ji:2013dva,Ji:2014hxa,Ji:2020ect}. So far, LaMET has been widely used in calculating quark isovector distribution functions~\cite{Lin:2014zya,Alexandrou:2015rja,Chen:2016utp,Alexandrou:2016jqi,Alexandrou:2018pbm,Chen:2018xof,Lin:2018pvv,Liu:2018uuj,Alexandrou:2018eet,Liu:2018hxv,Chen:2018fwa,Izubuchi:2019lyk,Shugert:2020tgq,Chai:2020nxw,Lin:2020ssv,Fan:2020nzz}, generalized parton
distributions~\cite{Chen:2019lcm,Alexandrou:2019dax}, distribution amplitudes (DAs)~\cite{Zhang:2017bzy,Chen:2017gck,Zhang:2020gaj}, and transverse-momentum-dependent distributions~\cite{Shanahan:2019zcq,Shanahan:2020zxr,Zhang:2020dbb}. 

LaMET suggests to calculate time-independent physical distributions in a finite momentum nucleon, which
are Euclidean observables. Such finite-momentum quantities can then be matched to light-front (LF)
parton properties using effective field theory techniques~\cite{Ji:2020byp}. For example, for the collinear quark distributions, it has been suggested to first compute a Euclidean space correlation (or quasi-light-front (quasi-LF) correlation) on the lattice, 
\bea\label{eq:quasi}
\tilde{h}(z, P_z) = \langle P| O_{\gamma_t}(z)|P \rangle,
\eea
where $|P\rangle$ is the nucleon state with a large momentum $P$ and the non-local operator is
\bea\label{eq:operator}
O_{\Gamma}(z)=\bar{\psi}(0) \Gamma U(0,z) \psi (z),
\eea
where $\psi$, $\bar{\psi}$ denote the bare quark field, $\Gamma$ is a Dirac structure,  and $U(0, z) = \exp(-ig\int_0^{z} dz' A_z(z'))$ is the Wilson link along the direction $z$, where $A^\mu$ is the gluon gauge potential. After renormalizing $\tilde{h}$ properly and matching it to some continuum scheme such as
dimensional regularization and (modified) minimal subtraction ($\overline{\rm MS}$), one obtains the so-called quark quasi-PDF via a Fourier transform, from which the quark PDF can be extracted through perturbative QCD matching. 

Renormalizing the quasi-LF correlation $\tilde{h}$ under lattice regularization is a non-trivial task. To see this, let us take, as a simple example, the matrix element of $O_\Gamma(z)$ in an off-shell quark state $|q\rangle$ with large Euclidean momentum $p^2$ (actually an off-shell truncated Green's function). It has the following form at 1-loop level~\cite{Alexandrou:2017huk},
\bea\label{eq:1-loop}
\langle O_{\Gamma}(z)\rangle =\Gamma\left(1+\gamma g^2 \textrm{log}(z^2/a^2)-m_{-1}\frac{z}{a}+……\right),\nonumber\\
\eea
where $g$ is the bare gauge coupling and $a$ is the lattice spacing. Unlike the coefficient of the logarithm $\gamma$, the coefficient of the linear divergence $m_{-1}$ can be sensitive to the details of the fermion and gauge actions on the lattice as those actions themselves define the lattice regularization. The linear-divergence term is proportional to the Wilson link length $z/a$ in lattice units, and it can be exponentially large at large $z$ or small $a$, when higher-order corrections are included. Thus, the linear divergence effect has to be removed before the lattice calculated quasi-PDF can be extrapolated to the continuum limit $a\rightarrow0$ and eventually matched to the continuum-scheme PDF.

Theoretical studies so far have shown that the quasi-PDF operator is multiplicatively renormalizable~\cite{Ji:2015jwa, Ji:2017oey,Ishikawa:2017faj,Green:2017xeu} in a continuum theory. On the lattice, due to non-commutativity of the limit $z\to 0$ and $a\to 0$, it has been suggested recently~\cite{Ji:2020brr} to use
a hybrid scheme to renormalize
the large and short distance correlations separately, which has the advantage of avoiding certain discretization effects and undesired infrared effects introduced in the renormalization stage. At small $z$ where
a short distance expansion is valid, one can use various ratio schemes, such as dividing by an off-shell quark matrix element of the quasi-PDF operator in the regularization-independent momentum subtraction (RI/MOM) scheme~\cite{Green:2017xeu,Chen:2017mzz,Alexandrou:2017huk} or by a hadron matrix element in the rest frame $\tilde{h}(z, P_z=0)$~\cite{Orginos:2017kos,Izubuchi:2018srq}. At large distance, one can
directly remove the Wilson line self-energy effect by subtracting the
power divergence~\cite{Chen:2016fxx,Ji:2017oey,Ishikawa:2017faj,Green:2017xeu}. Apart from using the RI/MOM factor or the rest-frame 
hadron matrix element, other methods have been suggested to extract the 
linear divergence, including Wilson loop~\cite{Chen:2016fxx,Zhang:2017bzy}, 
vacuum expectation values $\langle O_{\Gamma}(z)\rangle $~\cite{Braun:2018brg}, and gauge fixed Wilson link~\cite{Ji:2020ect}, and so on~\cite{Ji:2020brr}. 

However, numerically subtracting the linear divergence is an extremely delicate exercise. First, linear divergence could be sensitive to computational systematics in lattice calculations. In data, there may be slight differences between two matrix elements with the same linear divergence. These differences may lead to a failure of the suggested methods, especially for small lattice spacings. Second, some of the renormalization factors have other problems. For example, the leading contribution to the vacuum expectation value of $\langle O_{\Gamma}(z)\rangle $ at short distance vanishes and therefore, it is numerically very challenging to obtain the relevant linear divergence. Finally, as we shall see, linearly-divergent chiral symmetry breaking effects for Wilson fermions may render the linear divergence non-universal~\cite{Zhang:2020rsx}. 

To understand better the effects of the linear divergence in LaMET applications, we study systematically the linear divergences of matrix elements for sets of lattice data, generated with different lattice actions. We propose a self-renormalization method to eliminate all divergences and discretization errors when data for several different lattice spacings are available. The idea of this method is to extract the renormalization factor and the residual contribution directly from the matrix element we want to renormalize, without using an additional matrix element for renormalization. We then match the empirically renormalized 
matrix elements to those in the continuum. Our method is largely model-independent, and each term of our fitting functions is motivated by physics considerations. Tests show that our method works well for all data sets considered, which include matrix elements calculated with the lattice spacings from $0.03$ fm to $0.12$ fm using the valence clover and overlap actions on the MILC~\cite{Bazavov:2012xda} $N_f=2+1+1$  and RBC~\cite{Blum:2014tka} $N_f=2+1$ configurations. We also find that the method works for matrix elements after applying smearing which is needed to improve the statistical precision.



The rest of the paper is organized as follows. We present the theory of perturbative renormalization in Sec.~\ref{sec:theory}. The definition of the matrix elements we calculate and the simulation setup can be found in Sec.~\ref{sec:setup}. In Sec.~\ref{sec:test} the linear divergence is analysed for the different ensembles. Sec.~\ref{sec:StrtoRe} presents our strategy to self-renormalize a matrix element and collects the results for the $O_{\gamma_t}(z)$ matrix elements in different states and calculated with different valence fermion actions. Sec.~\ref{sec:highorder} discusses other fitting options, which mainly differ in the treatment of higher-order terms. In Sec.~\ref{sec:test},~\ref{sec:StrtoRe}, and~\ref{sec:highorder}, we only analyze matrix elements without HYP smearing. In Sec.~\ref{sec:smearing}, we extend our method to HYP smeared cases. In Sec.~\ref{sec:otherresult}, we test the linear divergence in some vacuum state matrix elements, including Wilson loop, quasi-PDF operator, and gauge fixed Wilson link, for the HYP smeared cases. In Sec.~\ref{sec:summary}, we summarize the results and discuss issues to be addressed in further studies.
